\chapter{Problem Formulation}\label{ch:problem_formulation}

\section{Power-Minimizing Joint Design}\label{sec:joint_problem}

The system model in Chapter~\ref{ch:system_model} separates globally
cooperative communication transmission from target-specific sensing
participation.  The design objective is to minimize the total transmit power
while meeting robust communication, detection, tracking, AP-power, and
sensing-cluster requirements.  The original mixed-integer problem is
\begin{subequations}\label{prob:p1}
\begin{align}
 \underset{\{\vect{w}_k\},\{\mats{S}_p\},\{b_{mp}\}}{\operatorname{minimize}}
 \quad & \sum_{k\in\mathcal{K}}\|\vect{w}_k\|_2^2+
 \sum_{p\in\mathcal{P}}\Tr(\mats{S}_p) \tag{P1}\\
 \text{subject to}\quad
 & \eqref{eq:robust_sinr_system}, \quad \forall k\in\mathcal{K},
 \tag{P1-C1}\\
 & \eqref{eq:sensing_sinr_model}, \quad \forall p\in\mathcal{P},
 \tag{P1-C2}\\
 & \eqref{eq:pcrb_system}, \quad \forall p\in\mathcal{P},
 \tag{P1-C3}\\
 & \Tr(\mats{E}_m\mats{S}_p)\leq P_{\max}b_{mp},
 \quad \forall m\in\mathcal{M},\ p\in\mathcal{P}, \tag{P1-C4a}\\
 & \Tr(\mats{E}_m\mats{R}_X)\leq P_{\max},
 \quad \forall m\in\mathcal{M}, \tag{P1-C4b}\\
 & \sum_{m\in\mathcal{M}}b_{mp}=N_{\mathrm{req}},
 \quad \forall p\in\mathcal{P}^{\mathrm{active}}, \tag{P1-C5}\\
 & \mats{S}_p\succeq\mats{0},
 \quad \forall p\in\mathcal{P}, \tag{P1-C6}\\
 & b_{mp}\in\{0,1\},
 \quad \forall m\in\mathcal{M},\ p\in\mathcal{P}. \tag{P1-C7}
\end{align}
\end{subequations}
The total covariance $\mats{R}_X$ is given in \eqref{eq:total_covariance}.
Constraint (P1-C4a) controls only the dedicated sensing covariance of target
$p$, whereas (P1-C4b) is the physical AP transmit-power limit.  This
two-sided distinction is essential: replacing them by a single sum Big-M
constraint would describe a different system and could incorrectly gate
communication beamforming.

\section{Sources of Difficulty}\label{sec:problem_difficulties}

Problem (P1) is a mixed-integer non-convex program.  Its principal sources of
difficulty are as follows.
\begin{enumerate}
    \item The robust communication constraint contains a fractional quadratic
    SINR and must hold for infinitely many channel errors in the ball
    \eqref{eq:csi_uncertainty}.
    \item The matrix inverse in the PCRB constraint is nonlinear, even though
    its epigraph has an exact semidefinite representation.
    \item The binary sensing-cluster indicators create combinatorial
    complexity and couple topology to the target-specific sensing power gate.
    \item After covariance lifting, the communication covariances must be
    rank one to correspond to physical beamforming vectors.
\end{enumerate}
The sensing SINR, the dedicated sensing gate, the AP hardware constraint, and
the cluster-cardinality equality are affine once the covariance variables are
used.  The following chapter converts the robust and PCRB constraints into
linear matrix inequalities (LMIs), relaxes the two discrete/non-convex
structures in a controlled way, and recovers a physically validated binary
solution.

\section{Interpretation of the Resource Trade-Off}\label{sec:resource_tradeoff}

The association variables do not directly constrain user beamforming.
Instead, they decide which AP blocks may contribute to each target-specific
covariance $\mats{S}_p$.  Those sensing covariances both consume the common
per-AP power budget and, without receiver-side cancellation, contribute to
communication interference.  Hence a sensing topology that is favorable only
in geometric distance can still be physically infeasible after robust user
QoS and PCRB requirements are imposed.  This motivates the joint
topology-and-beamforming recovery procedure introduced in
Chapter~\ref{ch:solution_method}.
